\section{Human Feedback Data for Speech Naturalness} \label{sec:human_data}
\begin{table*}[!t]
\centering
\footnotesize
\renewcommand{\arraystretch}{1}  % add vertical breathing room
\setlength{\tabcolsep}{9pt}
\captionsetup{font=small}

\caption{\textbf{Annotation Rubric for Speech Naturalness.}
Each dialog is evaluated along multiple sub-dimensions that capture different aspects of conversational speech quality.}
\vspace{-0.6em}

\begin{tabular}{p{3.2cm} p{12.2cm}}
\toprule
\textbf{Sub-metric} & \textbf{Description} \\
\midrule
\textbf{Expressive Intensity} &
The degree of expressiveness, reflecting how emotionally vivid or subdued the speech sounds. \\

\textbf{Expressive Correctness} &
The appropriateness of the expressiveness w.r.t. the conversational context and semantic content. \\

\textbf{Intonation} &
The quality of intonation, including pitch variation, stress patterns, and overall prosodic naturalness. \\

\textbf{NSVS}&
The naturalness of non-speech vocalizations and filler words (e.g., breaths, pauses, hesitations). \\

\textbf{Mispronunciation} &
Whether mispronunciations are present in the speech, reflecting pronunciation accuracy. \\

\textbf{Pacing} &
The naturalness of the pacing or speaking rate, including timing, pauses, and temporal flow. \\

\midrule
\textbf{Human-likeness} &
An overall rating of how closely the speech resembles natural human conversational speech. \\
\bottomrule
\end{tabular}
\label{table:rubric}
\vspace{-1em}
\end{table*}


To reliably conduct speech naturalness assessment, we collect high-quality human annotations that capture both fine-grained attributes and overall human-likeness of conversational speech. These annotations serve as the foundation for training speech naturalness predictor.

\textbf{Data Annotation Rubric.}
To assess speech naturalness, we design a structured annotation rubric that decomposes naturalness into multiple relevant sub-metrics together with an overall \emph{human-likeness} score. Human raters evaluate each system response independently. All sub-metrics are rated on a 5-point Likert scale, where higher values indicate better quality. One exception is for mispronunciation, which is annotated using a categorical scale: 1 means a mispronunciation is present, 2 represents no mispronunciation, 3 means not sure. The details are provided in Table~\ref{table:rubric}.

\textbf{Conversational TTS Dataset (ConvTTS).}  
ConvTTS dataset is collected using an in-house TTS system capable of producing two-channel dialogues. The model generates synthetic dialogue between a user and a system by conditioning on reference speech examples along with textual instructions. The resulting datasets contains 6,579 dialogue samples, each of length 49.4 $\pm$19.3 seconds. We partition the samples into 4,579 training samples, 1,000 development samples, and 1,000 test samples, corresponding to 62.8, 13.7, and 13.8 hours of audio, respectively. Each sample is annotated for audio naturalness by human evaluators at the dialogue level. A subset of the samples also include an additional rating of the assistant speech reponse, totaling to 3,263 training, 711 development, and 717 test samples. On average, each audio recording is labeled by 6.8 annotators, with the majority of samples evaluated by at least five different human raters.

\textbf{Full-duplex Conversational Dataset (FDX-Conv).} 
FDX-Conv dataset is collected from an in-house speech LLM designed for full-duplex interaction, i.e., the system can listen and speak at the same time, allowing overlapping user and system speech rather than strictly turn-based exchanges. The dataset contains 490 dialogue samples, each with an averaged duration of 41.1 $\pm$ 19.0 seconds. Similarly, each sample has been rated by no fewer than five human raters. Since the speech is generated by a model distinct from the ConvTTS system, we treat FDX-Conv as an out-of-domain (OOD) benchmark for assessing the generalization capability of naturalness predictors.


\textbf{Inter-Rater Consistency of Human-Likeness Ratings.}
To ensure the reliability of the collected annotations, we measure inter-rater consistency for overall human-likeness ratings. Specifically, for each audio sample, we randomly partition the set of annotators into two disjoint subsets, compute the average rating within each subset, and measure the Pearson correlation coefficient between the two averages. The Pearson correlation coefficients are $0.533$ for ConvTTS and $0.532$ for FDX-Conv (see Figure~\ref{fig:human_pcc} in Appendix), indicating strong inter-rater agreement. These results also serve as an empirical upper bound on the consistency between an automatic naturalness predictor and human raters. We further analyze how individual sub-metrics relate to the overall human-likeness score in Appendix~\ref{sec:data_analysis}.